% 球面 S² 参数化 (Sphere S^2 Parametrization)
% 描述: 展示二维球面的多种参数化方法
% 数学概念: 球面、球坐标、立体投影、黎曼球面
\documentclass[border=10pt]{standalone}
\usepackage{tikz}
\usepackage{amsmath,amssymb}
\usetikzlibrary{arrows.meta,positioning,calc,3d,patterns,decorations.pathmorphing}

\begin{document}
\begin{tikzpicture}[
    >=Stealth,
    scale=1.0
]

% ============= 左侧: 三维球面 =============
\begin{scope}[xshift=-5cm, yshift=0.5cm]
    % 标题
    \node[font=\bfseries] at (0, 3) {三维球面 $S^2$};
    
    % 球面轮廓
    \draw[color=blue!60!black, line width=1pt, fill=blue!10, opacity=0.6] 
        (0, 0) circle (2);
    
    % 赤道
    \draw[color=red!70!black, line width=1pt] 
        (0, 0) ellipse (2 and 0.6);
    
    % 经线 (前半)
    \draw[color=green!60!black, line width=0.8pt] 
        (0, -2) arc (-90:90:0.4 and 2);
    \draw[color=green!60!black, line width=0.8pt, dashed] 
        (0, 2) arc (90:270:0.4 and 2);
    
    % 北纬线
    \draw[color=orange!80!black, line width=0.8pt, dashed] 
        (0, 1) ellipse (1.73 and 0.52);
    
    % 坐标轴
    \draw[->, line width=0.8pt, color=black!60] (0, 0) -- (2.5, 0) node[right] {$x$};
    \draw[->, line width=0.8pt, color=black!60] (0, 0) -- (0, 2.5) node[above] {$z$};
    \draw[->, line width=0.8pt, color=black!60] (0, 0) -- (-0.8, -0.8) node[below left] {$y$};
    
    % 点P标注
    \coordinate (P) at (1.2, 1);
    \fill[black] (P) circle (2pt);
    \node[above right, font=\small] at (P) {$P$};
    
    % 角度标注
    \draw[->, color=purple!70!black, line width=0.8pt] (0.8, 0) arc (0:40:0.8);
    \node[color=purple!70!black, font=\small] at (1, 0.3) {$\theta$};
    
    % 半径
    \draw[dashed, color=gray] (0, 0) -- (P);
    \node[color=gray, font=\tiny, above left] at (0.6, 0.5) {$r=1$};
    
    % 北极标记
    \fill[black] (0, 2) circle (2pt);
    \node[above] at (0, 2) {$N$};
    \fill[black] (0, -2) circle (2pt);
    \node[below] at (0, -2) {$S$};
\end{scope}

% ============= 中间: 球坐标参数化 =============
\begin{scope}[xshift=1cm, yshift=0.5cm]
    % 标题
    \node[font=\bfseries] at (0, 3) {球坐标参数化};
    
    % 参数公式框
    \draw[rounded corners=5pt, fill=yellow!15, inner sep=10pt] 
        (-2.2, -1) rectangle (2.2, 2);
    
    \node[anchor=north, font=\small, align=center] at (0, 1.8) {
        $\begin{cases}
            x = \sin\theta \cos\phi \\
            y = \sin\theta \sin\phi \\
            z = \cos\theta
        \end{cases}$\\[8pt]
        $\theta \in [0, \pi]$ \\[3pt]
        $\phi \in [0, 2\pi)$
    };
    
    % 角度说明
    \node[anchor=north, font=\scriptsize, text width=4.5cm, align=left] 
        at (0, -1.5) {
        $\theta$: 极角(与z轴夹角)\\
        $\phi$: 方位角(xy平面)
    };
    
    % 坐标卡说明
    \node[anchor=north, font=\scriptsize, text width=4.5cm, align=center, 
          draw=green!50, fill=green!5, rounded corners, inner sep=5pt] 
        at (0, -3) {
        需要两个坐标卡覆盖\\
        由球极投影给出
    };
\end{scope}

% ============= 右侧: 球极投影 =============
\begin{scope}[xshift=5.5cm, yshift=0.5cm]
    % 标题
    \node[font=\bfseries] at (0, 3) {球极投影};
    
    % 球面简化表示
    \draw[color=blue!60!black, line width=1pt, fill=blue!10, opacity=0.6] 
        (0, 0) circle (1.5);
    
    % 赤道平面
    \fill[gray!20, opacity=0.5] (-2.5, -1.5) rectangle (2.5, -2.5);
    \draw[color=gray!60, line width=0.8pt] (-2.5, -1.5) -- (2.5, -1.5);
    \node[font=\scriptsize, color=gray!60] at (2.2, -1.8) {$\mathbb{C} \cong \mathbb{R}^2$};
    
    % 北极
    \fill[red!70!black] (0, 1.5) circle (3pt);
    \node[above, color=red!70!black] at (0, 1.5) {$N$};
    
    % 球面上一点
    \coordinate (Q) at (0.9, 0.5);
    \fill[black] (Q) circle (2pt);
    \node[above right, font=\small] at (Q) {$z$};
    
    % 投影点
    \coordinate (Qz) at (1.8, -1.5);
    \fill[black] (Qz) circle (2pt);
    \node[below, font=\small] at (Qz) {$\pi(z)$};
    
    % 投影线
    \draw[color=orange!80!black, line width=1pt, dashed] (0, 1.5) -- (Qz);
    
    % 公式
    \node[anchor=north, font=\scriptsize] at (0, -3) {
        $\pi: S^2 \setminus \{N\} \to \mathbb{R}^2$\\[3pt]
        $\pi(x, y, z) = \left(\frac{x}{1-z}, \frac{y}{1-z}\right)$
    };
    
    % 逆映射
    \node[anchor=north, font=\tiny, text width=5cm, align=center] 
        at (0, -4.3) {
        逆: $\pi^{-1}(u, v) = \frac{(2u, 2v, u^2+v^2-1)}{u^2+v^2+1}$
    };
\end{scope}

% ============= 底部: 数学性质 =============
\node[anchor=north, font=\small, draw=blue!30, fill=blue!5, rounded corners, inner sep=10pt] 
    at (4, -5) {
    \begin{tabular}{ll}
        \textbf{代数拓扑不变量:} & \\
        基本群: & $\pi_1(S^2) = \{1\}$ (单连通) \\
        同调群: & $H_0 = \mathbb{Z}, \; H_1 = 0, \; H_2 = \mathbb{Z}$ \\
        Euler示性数: & $\chi(S^2) = 2$ \\
        可定向性: & 可定向 \\
        切丛: & $TS^2$ 非平凡 (毛球定理) \\
        分类: & 亏格0闭曲面
    \end{tabular}
};

\end{tikzpicture}
\end{document}
